\section{Conclusion}
\label{sec:conclusion}

This paper presented a hybrid architecture for prioritized multi-AGV scheduling in a multi-floor smart-manufacturing layout: a central server assigns jobs, splits routes into floor legs at lifts, and books shared lift trips; floor servers plan local route-leg schedules; lift-lobby waiting points absorb delays near floor-level conflicts and lifts; an independent checker validates every complete plan. On the tested three-floor layout (27 nodes, 4 lifts), ten job sets of 100 jobs and 20 AGVs, the hybrid configuration reduced mean makespan by \RedMakespanCAabs\%, tardiness by \RedTardinessCAabs\%, and imbalance by \RedImbalanceCAabs\% relative to the paper-style centralized configuration, with busiest-server CPU down \RedCpuMaxCAabs\%. B2 and C matched schedule outcomes on every seed, so measured gains come from lobby waiting and FCFS leg ordering; the floor decomposition changes where reservation state and computation sit.

The reduction in tardiness is substantial, but mean total tardiness remains \TardinessC~s per 100-job seed for C. Conflict resolution and decomposition alone do not ensure deliveries meet their targets: assignment is greedy and priority-dependent, and delay accumulated through shared routes, lift use, and safety separation can still push deliveries past desired times. Reducing residual tardiness requires changes to assignment, sequencing, and resource coordination beyond distributing the planning workload. These results apply to the tested layout, fleet sizes, seeds, and parameters; they establish neither optimality nor a general scalability claim.

\paragraph{Future work.} Dynamic job arrivals, battery limits and charging, multiple lifts per floor with explicit capacity, decentralized lift coordination, and comparisons with auction-based assignment, DRL, and CBS on common instances with shared verification rules.
